\chapter{PROBLEMÁTICA}
\label{cap:problematica}

\section{DESCRIPCIÓN DE LA PROBLEMÁTICA}

Los censos de árboles urbanos constituyen una actividad de gestión ambiental de
primera importancia para los municipios, ya que permiten caracterizar el patrimonio
verde de la ciudad, evaluar su cobertura y planificar intervenciones de mantenimiento.
Esta información es esencial para mitigar riesgos ambientales, prevenir accidentes
por caídas de árboles y optimizar recursos municipales.

Arbocensus, desarrollado en la Universidad de los Andes en el marco del proyecto
Fondef ID21110360, censa el arbolado urbano mediante visión artificial y
participación ciudadana \citep{binder-cortes2023,pizarro2024}. Las campañas
ejecutadas hasta ahora produjeron el catastro de árboles georreferenciados que esta
iteración utiliza como punto de partida (la Sección~\ref{sec:alcance-desarrollo}
detalla las iteraciones que lo produjeron); ese catastro es, sin embargo, una fotografía
única de cada árbol.

Para mejorar la calidad de los datos en futuras campañas de monitoreo, así como
para entrenar modelos de inteligencia artificial, surge una necesidad crítica:
realizar un re-censo estructurado que permita construir un historial de
observaciones por árbol ---objetivo declarado del proyecto \citep{pizarro2024}---
y provea datos etiquetados que permitan clasificar los árboles por nivel de riesgo
(riesgoso/no-riesgoso). Con este fin resulta necesario repartir el arbolado entre
los equipos de trabajo mediante rutas eficientes de desplazamiento.

La escala real del problema se evidencia en las dos bases de datos que dejaron las
iteraciones anteriores del proyecto \citep{binder-cortes2023,pizarro2024}. La
más antigua, la base \texttt{legacy\_app} de la aplicación de censado previa desplegada
en Heroku, registra 17\,208 muestras de censo asociadas por código QR, que corresponden a
13\,081 árboles distintos. De ellos, los validados por el canal de medición automática
(\texttt{mltree}) y con una posición geográfica única ---determinada mediante la mediana
de las muestras de cada QR--- suman 1\,178 árboles importables.
La posterior, la
base de datos \texttt{legacy\_api}, contiene 429 árboles georreferenciados, 3\,961
muestras (\textit{samples}), 10\,008 fotografías y 48 áreas de censo distribuidas en 10
campañas.
Al combinar ambas fuentes, el conjunto de árboles a rutear asciende a 1\,607 nodos,
un tamaño que sitúa el problema fuera del alcance de la planificación manual y que
impone restricciones sobre la construcción de la matriz de costos (RE-01, RE-02).

Actualmente, con las coordenadas geográficas disponibles de los árboles, la organización
de los censadores en terreno representa un desafío logístico significativo: dado un
conjunto de $n$ puntos geográficos correspondientes a árboles distribuidos en el espacio
urbano, se requiere generar un conjunto de rutas que garanticen la cobertura total,
minimicen tiempos de desplazamiento y respeten jornadas acotadas temporalmente.

La planificación inadecuada de rutas genera consecuencias operacionales inmediatas:
viajes innecesarios entre puntos lejanos, sobrecarga de ciertos equipos mientras otros
subutilizan su tiempo, imposibilidad de cubrir todos los árboles en una jornada y
aumento de costos logísticos.

\section{ESTADO ACTUAL DEL PROBLEMA}
\label{sec:estado-actual}

Ante la ausencia de herramientas especializadas de optimización, la planificación de
rutas para censos urbanos queda limitada a métodos manuales ---hojas de cálculo o
herramientas de sistemas de información geográfica (GIS) de propósito general--- sin capacidad de optimización combinatoria. Los administradores de bases de datos arbóreas deben
asignar manualmente árboles a censadores con apenas factores espaciales rudimentarios
(proximidad euclidiana), sin considerar costos reales de desplazamiento ni
restricciones de jornada.

Las soluciones \textit{greedy} (vecino más cercano, \textit{clustering} geométrico) generan resultados
subóptimos: viajes innecesariamente largos, desbalance severo en el número de árboles
asignado a cada censador o incapacidad para cumplir restricciones temporales por jornada
laboral. La comparación del Capítulo~\ref{cap:resultados} ilustra las dos primeras
deficiencias sobre el censo real: el heurístico del vecino más cercano acumula más
desplazamiento total que el modelo de optimización y produce una ruta degenerada que
dedica más de dos horas y media de caminata a cinco árboles, además de saturar cada
jornada al borde de su cota, sin holgura ante imprevistos de terreno.

\section{ALCANCE}
\label{sec:alcance-desarrollo}

La plataforma de censado se construyó en iteraciones sucesivas desde 2021: una
primera versión web orientada a dispositivos móviles que permitió tomar muestras de
árboles ---mediante lectura de códigos QR y captura de GPS--- organizadas en
campañas, zonas y equipos; un rediseño que incorporó elementos de ludificación
---puntos, logros, desafíos y tablas de posiciones--- para incentivar la
participación \citep{binder-cortes2023}; y una tercera iteración que integró el
procesamiento de imágenes del Sistema de Análisis de Imágenes de Árboles (SAIdA) y
un modelo de estimación de cobertura del censo \citep{pizarro2024}.
% Historia verificada contra los repositorios de código del proyecto (jul 2026):
%   - github.com/mrecabarren/arbocensus: commit inicial 2021-10-01 (Matías Recabarren,
%     profesor guía). Django + templates HTML/JS, escaneo de QR con cámara, Google auth,
%     despliegue en Heroku, modelos Sampling/{sample,question,answer,step,results}.
%     = primera versión: app web de censado orientada a móviles. Fija el inicio en 2021.
%   - github.com/Arbocensus/ArbocensusGamification: primer commit de código 2023-03-17
%     (Julio Cortés), repo de la organización creado 2023-10-25. fastAPI + alembic,
%     modelos badge/campaign/challenge/achievement/score/ranking = capa de ludificación.
%     => la ludificación de Binder y Cortés se desarrolló en 2023; la memoria se cita como
%     2022, pero se asume que es un error debido a las fechas de los commits.
% ASUMIDO (no confirmable con las fuentes disponibles):
%   - Autoría de la primera versión: Binder y Cortés la atribuyen a un trabajo de título
%     anterior (Castro y Díaz), pero el historial git muestra como autores a Recabarren,
%     "Richard Le May" y "Track-tor", no esos nombres. No se afirma autoría en el texto.
%   - Discrepancia de año de inicio entre fuentes (Binder/Cortés: 2022; Pizarro:
%     2021) queda resuelta a 2021 por el commit inicial del repositorio original.

Ninguna de esas iteraciones abordó la planificación de rutas: la asignación de
árboles a censadores y el orden de visita se decidieron manualmente en cada campaña.
Esta iteración agrega esa pieza faltante: un motor de optimización que determina
automáticamente el número de equipos necesarios, asigna los árboles entre ellos y
ordena las visitas sobre la red vial real, además del ciclo operacional que lo rodea
---importación del arbolado, publicación del plan, portal móvil del censador y
seguimiento del avance---, que el Capítulo~\ref{cap:solucion} detalla.
